\section*{Bab \ref{ch:g7:atoms-and-molecules} --- Atom dan Molekul}

\begin{solution}{exo:g7:atoms-and-molecules:1}
\omterm{def:g7:atoms-and-molecules:atom}{Atom} adalah salah satu partikel yang sangat kecil yang menyusun semua
materi (ada sekitar seratus jenis). \omterm{def:g7:atoms-and-molecules:molecule}{Molekul} adalah sekelompok \omterm{def:g7:atoms-and-molecules:atom}{atom} yang
terikat erat, bagian terkecil sebuah zat yang masih merupakan zat itu.
\end{solution}

\begin{solution}{exo:g7:atoms-and-molecules:2}
\ce{H}, \ce{C}, \ce{N}, \ce{O}, \ce{Na}, \ce{Fe}.
\end{solution}

\begin{solution}{exo:g7:atoms-and-molecules:3}
1 \omterm{def:g7:atoms-and-molecules:atom}{atom} karbon dan 2 \omterm{def:g7:atoms-and-molecules:atom}{atom} oksigen: seluruhnya 3 \omterm{def:g7:atoms-and-molecules:atom}{atom}.
\end{solution}

\begin{solution}{exo:g7:atoms-and-molecules:4}
Air, oksigen, metana, nitrogen.
\end{solution}

\begin{solution}{exo:g7:atoms-and-molecules:5}
\ce{NO2}.
\end{solution}

\begin{solution}{exo:g7:atoms-and-molecules:6}
\ce{O2} adalah satu \omterm{def:g7:atoms-and-molecules:molecule}{molekul} yang tersusun atas dua \omterm{def:g7:atoms-and-molecules:atom}{atom} oksigen yang
bergabung; \ce{2O} adalah dua \omterm{def:g7:atoms-and-molecules:atom}{atom} oksigen yang terpisah. \ce{CO} adalah
\omterm{def:g7:atoms-and-molecules:molecule}{molekul} yang terdiri atas satu \omterm{def:g7:atoms-and-molecules:atom}{atom} karbon dan satu \omterm{def:g7:atoms-and-molecules:atom}{atom} oksigen (dua
huruf kapital, dua \omterm{def:g7:atoms-and-molecules:symbol}{lambang}); \ce{Co} adalah satu \omterm{def:g7:atoms-and-molecules:atom}{atom} kobalt (satu
\omterm{def:g7:atoms-and-molecules:symbol}{lambang}, huruf kapital lalu huruf kecil).
\end{solution}

\begin{solution}{exo:g7:atoms-and-molecules:7}
\ce{3H2O}: $3\times 2 = 6$ \omterm{def:g7:atoms-and-molecules:atom}{atom} hidrogen dan $3\times 1 = 3$ \omterm{def:g7:atoms-and-molecules:atom}{atom}
oksigen. \ce{5O2}: tidak ada hidrogen, $5\times 2 = 10$ \omterm{def:g7:atoms-and-molecules:atom}{atom} oksigen.
\end{solution}

\begin{solution}{exo:g7:atoms-and-molecules:8}
Hitam adalah karbon, merah adalah oksigen: satu karbon, dua oksigen,
\ce{CO2}, karbon dioksida.
\end{solution}

\begin{solution}{exo:g7:atoms-and-molecules:9}
Air menurut Dalton adalah \ce{HO}. \omterm{def:g7:atoms-and-molecules:molecule}{Molekul} yang sebenarnya, \ce{H2O},
memiliki satu \omterm{def:g7:atoms-and-molecules:atom}{atom} hidrogen lagi: \omterm{def:g7:atoms-and-molecules:molecule}{molekul} Dalton kekurangan satu \omterm{def:g7:atoms-and-molecules:atom}{atom}.
\end{solution}

\begin{solution}{exo:g7:atoms-and-molecules:10}
Air murni bukan \omterm{def:g6:pure-substances-and-mixtures:mixture}{campuran}: semua \omterm{def:g7:atoms-and-molecules:molecule}{molekulnya} identik, yaitu \omterm{def:g7:atoms-and-molecules:molecule}{molekul}
\ce{H2O}. \omterm{def:g4:air-a-mixture-of-gases:air}{Udara} adalah \omterm{def:g6:pure-substances-and-mixtures:mixture}{campuran}: \omterm{def:g4:air-a-mixture-of-gases:air}{udara} mengandung \omterm{def:g7:atoms-and-molecules:molecule}{molekul} nitrogen,
\omterm{def:g7:atoms-and-molecules:molecule}{molekul} oksigen, \omterm{def:g7:atoms-and-molecules:atom}{atom} argon, dan sedikit \omterm{def:g7:atoms-and-molecules:molecule}{molekul} karbon dioksida.
\end{solution}

\begin{solution}{exo:g7:atoms-and-molecules:11}
Sepuluh juta \omterm{def:g7:atoms-and-molecules:atom}{atom} membentuk \qty{1}{mm}, jadi \qty{1}{cm} $= \qty{10}{mm}$
memerlukan $10 \times$ sepuluh juta $=$ seratus juta \omterm{def:g7:atoms-and-molecules:atom}{atom}.
$\qty{100}{m} = \num{100000}$~mm memerlukan $\num{100000}\times$
sepuluh juta $=$ sejuta juta \omterm{def:g7:atoms-and-molecules:atom}{atom}.
\end{solution}

\begin{solution}{exo:g7:atoms-and-molecules:12}
$6 + 12 + 6 = 24$ \omterm{def:g7:atoms-and-molecules:atom}{atom}. Hidrogen: $\frac{12}{24} = \qty{50}{\percent}$.
Karbon: $\frac{6}{24} = \qty{25}{\percent}$ (dan oksigen
\qty{25}{\percent} sisanya).
\end{solution}

\begin{solution}{pb:g7:atoms-and-molecules:1}
\textbf{1.} Nitrogen \ce{N2}, 2 \omterm{def:g7:atoms-and-molecules:atom}{atom}; oksigen \ce{O2}, 2 \omterm{def:g7:atoms-and-molecules:atom}{atom}; karbon
dioksida \ce{CO2}, 3 \omterm{def:g7:atoms-and-molecules:atom}{atom}.

\textbf{2.} \omterm{def:g7:atoms-and-molecules:atom}{Atom-atom} argon tidak bergabung menjadi \omterm{def:g7:atoms-and-molecules:molecule}{molekul}: argon
tersusun atas \omterm{def:g7:atoms-and-molecules:atom}{atom-atom} tunggal.

\textbf{3.} Nitrogen: $78 \times 2 = 156$ \omterm{def:g7:atoms-and-molecules:atom}{atom}. Oksigen: $21\times 2 =
42$ \omterm{def:g7:atoms-and-molecules:atom}{atom}.

\textbf{4.} $156 + 42 + 1 = 199$ \omterm{def:g7:atoms-and-molecules:atom}{atom}.

\textbf{5.} $\frac{156}{199} \approx 0.78$, sekitar \qty{78}{\percent}.

\textbf{6.} Di dalam \omterm{def:g7:atoms-and-molecules:molecule}{molekul} oksigen: $16\times 2 = 32$; di dalam
\omterm{def:g7:atoms-and-molecules:molecule}{molekul} karbon dioksida: $4 \times 2 = 8$; seluruhnya $32 + 8 = 40$.

\textbf{7.} $42 - 40 = 2$ \omterm{def:g7:atoms-and-molecules:atom}{atom} oksigen hilang: setara satu \omterm{def:g7:atoms-and-molecules:molecule}{molekul}
oksigen. \omterm{def:g7:atoms-and-molecules:atom}{Atom-atom} itu masuk ke dalam \omterm{def:g7:atoms-and-molecules:molecule}{molekul} air, \ce{H2O}, yang dibuat
tubuh dengan hidrogen dari makanannya; air itu keluar sebagian sebagai
uap air dalam embusan napas (dihilangkan dalam angka-angka ini).

\textbf{8.} Dihirup: 0. Diembuskan: 4 (satu dalam setiap \ce{CO2}).

\textbf{9.} Dari makanan yang dipakai tubuh: gula dan lemak mengandung
\omterm{def:g7:atoms-and-molecules:atom}{atom} karbon, dan tubuh melepaskannya dalam karbon dioksida yang
diembuskannya.

\textbf{10.} $21 - 16 = 5$ \omterm{def:g7:atoms-and-molecules:molecule}{molekul} oksigen hilang dan 4 \omterm{def:g7:atoms-and-molecules:molecule}{molekul} karbon
dioksida muncul.

\textbf{11.} Nitrogen: dua bola biru yang dihubungkan tiga tongkat.
Oksigen: dua bola merah yang dihubungkan dua tongkat. Karbon dioksida:
satu bola hitam di antara dua bola merah, masing-masing terhubung
dengannya oleh dua tongkat, dalam satu garis lurus. Argon: satu bola
biru pucat tunggal.

\textbf{12.} $156 + 32 + 12 + 1 = 201$ \omterm{def:g7:atoms-and-molecules:atom}{atom} (keempat \omterm{def:g7:atoms-and-molecules:molecule}{molekul} karbon
dioksida mengandung 12 \omterm{def:g7:atoms-and-molecules:atom}{atom}). Itu $201 - 199 = 2$ \omterm{def:g7:atoms-and-molecules:atom}{atom} lebih banyak
daripada \omterm{def:g4:air-a-mixture-of-gases:air}{udara} yang dihirup: 4 \omterm{def:g7:atoms-and-molecules:atom}{atom} karbon bertambah dari makanan,
dikurangi 2 \omterm{def:g7:atoms-and-molecules:atom}{atom} oksigen yang pergi dalam air.
\end{solution}
